\section{Conclusion}
\label{sec:conclusion}

We opened this paper with a near-zero number --- Pearson $|r|(\text{SE}_{N5}, \text{min\_logprob}) = 0.026$ --- and asked what it means. We can now close the loop: it means that Semantic Entropy and minimum log-probability are genuinely distinct uncertainty signals, each carrying information about LLM hallucination through mechanisms that are empirically as well as algebraically independent. SE asks whether a model is semantically consistent across stochastic samples; min\_logprob asks where a model is least confident in a single deterministic pass. The near-zero correlation is the empirical signature of these fundamentally different uncertainty operations, not a statistical coincidence.

In this work, we directly measured the statistical independence of $\text{SE}_{N5}$ and min\_logprob as uncertainty signals for hallucination detection in open-weight LLMs on short-answer factual QA. Our main contributions are: (1) Pearson $|r|=0.026$ confirming near-orthogonality at $N=2500$; (2) partial $R^2=0.0005$ (LRT $p=0.442$) showing that independence does not yield linear conditional predictive utility on this benchmark --- a powered null result for the pre-registered effect size, with the gate outcome (EXPLORE\_N10) disclosed transparently; (3) Spearman $\rho(\text{SE}, \text{judge})=-0.026$ validating cross-model LM-judge as a low-circularity evaluation protocol; and (4) a fully reproducible, checkpoint-aware pipeline that completed the $N=2500$ evaluation.

Future work includes nonlinear ensemble evaluation (h-e1-v2), extending to $N=10$ stochastic samples per prompt (per the gate recommendation), cross-model transfer characterization (Llama-calibrated ensemble on Qwen-2.5-7B TruthfulQA), orthogonality zone analysis, bootstrapped CIs on primary metrics, and seed-stability verification. Hallucination detection in open-weight LLMs is a prerequisite for safe deployment in high-stakes settings. We hope this work demonstrates that careful statistical characterization of signal independence --- not just standalone AUROC comparison --- is the right foundation for principled ensemble design.
